\section{Quotients and first-level ramification}\label{sec:quotients}
We record the formal-group facts needed to pass between truncation levels. The residue field in this section is algebraically closed and $p\ge5$.

\begin{lemma}\label{lem:tame-good}
A height-two group with potentially supersingular reduction acquires good formal reduction over a tame cyclic extension of degree dividing $N=p^2-1$. Its finite special character has eigenvalues $\zeta^c,\zeta^{pc}$, and its generic determinant is $\chi_p\epsilon$, with $\epsilon$ the finite norm character.
\end{lemma}
\begin{proof}
The endomorphism algebra of the special height-two Honda group is the division algebra
\[
 D=\Q_{p^2}\langle\Pi\rangle,
 \qquad \Pi^2=p,\qquad \Pi a=\sigma(a)\Pi
\]
for $a$ in the unramified coefficient field and its nontrivial Frobenius $\sigma$. Its maximal order has residue field $\F_{p^2}$ and pro-$p$ reduction kernel. An element of order $p$ would generate the field $\Q_p(\zeta_p)$ of degree $p-1>2$, whereas every element of $D$ satisfies a reduced quadratic polynomial. Thus every finite subgroup of its units has order prime to $p$, by Cauchy's theorem, and reduction embeds it into $\F_{p^2}^{\times}$.

The kernel of the special inertia action may be removed from a good extension: the rational representation then has trivial inertial type and is crystalline, and the integral Barsotti--Tate classification applies to the original lattice, not merely to its isogeny class. This is the crystalline descent and lattice statement in \citet[Theorem~2.2.7 and Proposition~A.6]{Kisin2006}; Tate's full faithfulness identifies the resulting model with the original good model after extension \citep[Theorem~4, p.~180, and Corollary~2, p.~181]{Tate1967}. The faithful special quotient therefore supplies the asserted tame extension. The two special eigenvalues are the unramified conjugates in $\F_{p^2}$, and taking reduced norms gives the determinant assertion. Two such groups have a common good extension of degree dividing $N$, since tame inertia is procyclic.
\end{proof}

We may enlarge the good field to degree $N$ and put $s=(p+1)e$. This field also kills the tame quotient of the residual representation. Indeed, a nontrivial residual wild group has a unique invariant line, so its tame quotient acts through the two diagonal $\F_p$-characters; if the wild group is trivial, a tame generator is semisimple with eigenvalues in $\F_{p^2}$. In both cases its order divides $N$. Whenever $p-1\mid e$, the integer $s$ is a multiple of $N$. In a two-slope polygon the parameter $B$ of~\eqref{eq:two-slope-B} satisfies
\[
           0<B<pe,\quad B\in\Z,\quad B\equiv c\pmod{p+1}.
\]
The integrality follows because the invariant high line has trivial character over $L$; the congruence follows because its tame character $\zeta^{c-B}$ is $\F_p$-valued. Over $L$, the first-level image is either trivial or the upper unipotent group of order $p$.

\begin{lemma}\label{lem:actual-quotients}
The quotient by the high line has parameter $pB$ if $B<e$ and $s-B$ if $B>e$; at $B=e$ its first-level polygon has one slope and every nonzero point has valuation $e$. If the first-level module is rational and $J$ is an inertia-stable low line, the quotient by $J$ has parameter $B/p$. A one-slope group has first-level value $e$, and its quotient by any stable line has parameter $e$, high value $pe$, and low value $e/p$.

These are good formal quotients over $L$ with the same finite special character. A comparison at level $p^n$ preserving the chosen lines induces a comparison of their quotients at level $p^{n-1}$.
\end{lemma}
\begin{proof}
The finite flat closure of a stable line $J$ defines an isogeny whose coordinate is
\[
                T\longmapsto U(T)\prod_{Q\in J}(T-T(Q)),
                 \qquad U(T)\in\Ocal_L[[T]]^\times.
\]
For the high line, every argument of value below $\alpha=s-B$ has its value multiplied by $p$. The image of the old first-level quotient has value $B$. If $B>e$, the inverse polygon over a nonzero high point has all $q$ roots of value $\alpha/q$; their quotient images give $q-p$ points of value $\alpha/p$, with $p-1$ old points of value $B$. This is parameter $s-B$. If $B<e$, the inverse roots have values $(s-pB)/p$ with multiplicity $p$, and $B/p$ with multiplicity $q-p$. Quotienting gives $p-1$ high points of value $s-pB$ and $q-p$ low points of value $B$, hence parameter $pB$. At equality the inverse polygon is a single segment and all quotient values are $e$.

For a low line, each nonzero old quotient class has a unique high representative. Its image has value $\alpha+(p-1)\beta=s-B/p$. The remaining images come from low inverse points of value $\beta/q$ and have value $\beta/p=B/p^2$. This gives parameter $B/p$. In the one-slope case every nonzero first-level point has value $e$, by the product formula. Each nonzero old quotient class has $p$ differences of value $e$, giving value $pe$; inverse points above the chosen line have value $e/q$ and quotient value $e/p$.

The special line is the Frobenius kernel, so the special quotient has tangent exponent $pc$ and the same unordered character $\{c,pc\}$. Tate full faithfulness extends the generic semilinear descent maps to the good models and preserves the finite flat closure, hence gives generic descent of the quotient. Finally, as actual finite Galois modules,
\[
 (M/J)[p^{n-1}]
       =\{P\in M[p^n]:p^{n-1}P\in J\}/J.
\]
This identity proves the precise truncation assertion.
\end{proof}

\begin{lemma}\label{lem:prime-residual-comparison}
In a two-slope comparison at level $p^2$ with nontrivial residual wild action, if at least one parameter $B_i$ is prime to $p$, then both parameters are equal and the special characters agree.
\end{lemma}
\begin{proof}
The common wild fixed line is the high line. The common first-level field $R/L$ has degree $p$. A low coordinate has value $B_i$ over $R$ and displacement $p(s-B_i)$. If $p\nmid B_1$, the exact prime-valuation formula gives
\[
                   b(R/L)=ps-(p+1)B_1.
\]
The unconditional inequality in the other coordinate gives $B_1\ge B_2$. If both are prime, symmetry gives equality. If $p\mid B_2$, equality is impossible. A common second-level lift of a low point has value $B_1/p^2$ over $R$; its denominator forces its point field to have degree $p^2$. Every first-level translation occurs, and the two normalized coordinates have values $B_1,B_2$. A low translation has displacement $p^2B_i$. The first model fixes its break as $(p^2-1)B_1$, whereas the second gives $p^2B_2\ge B_2+(p^2-1)B_1$, contradicting $B_1>B_2$. Equal parameters and the identified high-line character identify $c$ modulo $N$.
\end{proof}
